\chapter{Firefinger}
\sessioninfo{Session 10}{Mardi 10 novembre 2026}

Au lendemain de leur première nuit dans la jungle, les aventuriers reprirent leur descente vers le sud. Une routine s'installa rapidement : le chenal s'élargissait, le courant devenait plus paresseux, Azaka choisissait des passages absents de leur carte et de brèves pluies venaient régulièrement détremper la végétation. Sans leur guide, distinguer le cours principal des innombrables bras de la rivière aurait été presque impossible.

Leurs objectifs, au moins, restaient clairs. Ils devaient d'abord atteindre Firefinger et récupérer le masque d'Azaka, puis repartir à la recherche du Gardien Bouclier et de son amulette. S'ils découvraient en chemin l'origine du mal qui consumait Syndra, ce serait encore mieux.

\begin{narrateur}
Deux priorités supplémentaires s'ajoutaient à cette liste : ne pas mourir et trouver du cacao pour Morrÿs. Leur ordre d'importance faisait encore débat.
\end{narrateur}

Au soir du deuxième jour, Morrÿs installa un piège de chasse à proximité du campement et y déposa quelques restes de nourriture. Il espérait attirer l'un des petits animaux qui rôdaient autour de leurs provisions. Au matin, le mécanisme était toujours armé et l'appât intact. La jungle avait refusé de leur offrir le petit déjeuner.

Le troisième jour commença sous un ciel pâle et dans une humidité plus lourde encore qu'à l'ordinaire. Un brouillard épais recouvrait la rivière, dont l'eau était devenue trouble et agitée. Azaka reconnut les signes d'un orage tombé en amont : une crue sévère allait bientôt déferler dans le chenal. Elle leur déconseilla de remettre le canoë à l'eau.

Le groupe choisit malgré tout de poursuivre sa route à pied plutôt que d'attendre une journée entière. Obélik et Kael soulevèrent l'embarcation pendant qu'Azaka ouvrait le chemin le long de la rive. Devant eux, le grondement de la rivière enfla progressivement. L'eau monta jusqu'à quitter son lit et vint bientôt lécher leurs chevilles.

La brume devint si dense que Destos et Morrÿs perdirent les autres de vue. Le clerc ne distinguait plus rien à quelques pas, tandis que le nain apercevait encore vaguement leurs silhouettes devant eux. Ils se prirent par la main et accélérèrent avec prudence jusqu'à rejoindre le groupe. La pluie les rattrapa presque aussitôt et transforma le sol en un mélange de boue et d'eau.

\illustration[0.95\textwidth]{2026-11-10/brouillard.png}
{\textit{Dans l'épaisseur du brouillard, Morrÿs guida Destos pour rejoindre leurs compagnons.}}
{brouillard_firefinger}

Lorsque le niveau devint suffisant, ils retournèrent le canoë. Morrÿs s'y installa et Obélik entreprit de le tirer comme un traîneau, sous une pluie battante. L'embarcation resta d'abord obstinément collée dans la boue, puis céda d'un coup lorsque le goliath parvint à rompre son inertie. Tant qu'il ne s'arrêtait plus, il réussissait à la faire avancer.

\illustration[0.95\textwidth]{2026-11-10/crue.png}
{\textit{Alors que la rivière débordait, Obélik transforma le canoë en traîneau pour Morrÿs.}}
{crue_firefinger}

L'orage finit par s'éloigner et la rivière retrouva lentement un niveau moins inquiétant. À la fin de la journée, ils établirent leur camp à distance raisonnable du chenal. Morrÿs posa de nouveau son piège, tandis que Kael adapta son alarme afin que les ragondins et les plus petites créatures puissent franchir le périmètre sans réveiller tout le monde. La nuit s'écoula cette fois sans incident.

Au matin, un campagnol gisait dans le piège. L'animal n'avait pas survécu au mécanisme, mais fournirait au moins un peu de viande pour les repas suivants. Morrÿs considéra que l'expérience commençait enfin à porter ses fruits. Le groupe rangea le camp et remit le canoë à l'eau.

Azaka hésita un instant devant l'enchevêtrement des cours d'eau. La crue de la veille avait rendu le trajet difficile à reconnaître, mais Kael avait suivi leur progression avec suffisamment d'attention pour situer leur position. La guide retrouva rapidement ses repères et les ramena sur le bon chenal. Les quatrième et cinquième jours passèrent ensuite sans nouvelle rencontre.

\scenebreak

Au matin du sixième jour, la rivière se resserra encore. Le canoë s'engagea dans un cours d'eau peu profond, juste assez large pour leur permettre de progresser. La journée était claire et la lumière filtrait entre les branches lorsque, au détour d'un virage, un point orange apparut au-dessus de la jungle.

La lueur devint une flamme, puis une tour dressée au sommet d'un immense pilier rocheux. La vieille pierre, fissurée et couverte de lianes, dominait la rivière. Tout en haut, un brasier brûlait avec une vigueur constante sans que sa source soit visible. Autour du promontoire, de grandes silhouettes ailées tournoyaient, plongeaient entre les corniches ou se reposaient sur la roche.

Azaka cessa de pagayer.

\begin{quote}\itshape
Ça y est. C'est Firefinger. Mon masque est là-haut.
\end{quote}

\illustration[0.95\textwidth]{2026-11-10/08-promontoire-firefinger.png}
{\textit{Depuis la rivière, les aventuriers découvrirent Firefinger et les créatures qui tournoyaient autour de son sommet enflammé.}}
{decouverte_firefinger}

Le groupe observa les lieux depuis le canoë. Des échelles reliaient plusieurs ouvertures creusées dans la paroi, chacune paraissant mener à une cavité intérieure où il serait possible de reprendre pied. L'ascension se ferait donc par étapes d'une dizaine de mètres, avec autant d'occasions d'être repérés. Au pied du promontoire, un cadavre ancien reposait près de l'eau.

Morrÿs examina les ossements desséchés. Plusieurs fractures se concentraient sur un même côté du corps, sans trace évidente de morsure ni de griffure. Tout indiquait une chute depuis les hauteurs, même s'il était impossible de savoir si le corps avait atterri là ou si le courant l'y avait déposé. Azaka en conclut que cet inconnu avait probablement réussi à monter plus haut qu'elle.

La guide leur révéla n'avoir jamais atteint le premier palier. Lors de sa précédente tentative, les créatures l'avaient repérée presque immédiatement et elle avait dû renoncer plutôt que de rester seule, vulnérable sur une échelle. Elle les désigna comme des ptéranodons, des êtres natifs de Chult aussi prompts à chasser pour se nourrir que pour jouer. Ils pouvaient fondre sur leurs proies depuis le ciel, mais aussi lancer des projectiles avec leurs mains.

Attirer toute la colonie au pied de la tour fut rapidement écarté. Firefinger constituait leur bastion et les affronter tous reviendrait à se condamner. Le groupe envisagea de progresser palier après palier, chacun couvrant celui qui gravirait l'échelle suivante. Morrÿs proposa finalement d'attendre la nuit, puisque les créatures semblaient surtout actives pendant la journée.

\begin{narrateur}
Obélik disposait également d'une technique d'infiltration éprouvée : demeurer si parfaitement immobile qu'il en devenait invisible. Elle présentait malheureusement quelques limites lorsqu'il fallait escalader une tour.
\end{narrateur}

Ils éloignèrent le canoë et installèrent un campement de fortune, suffisamment loin pour ne pas attirer l'attention sans perdre Firefinger de vue. Pendant les dernières heures du jour, chacun prépara son équipement et observa les mouvements des ptéranodons. À la tombée de la nuit, les silhouettes ailées regagnèrent les corniches et le promontoire retrouva peu à peu son calme.

\scenebreak

La flamme du sommet et la lumière de la lune diffusaient une faible clarté autour de la tour. Kael et Azaka, capables de mieux distinguer les formes dans l'obscurité, encadrèrent la progression. Destos et Morrÿs avancèrent ensemble entre eux, tandis qu'Obélik fermait la marche. Tous s'efforçaient de poser chaque pied avec précaution.

Une branche sèche craqua sous la botte de Destos.

Là-haut, plusieurs ailes frémirent. Un ptéranodon quitta sa corniche et entama une ronde autour du promontoire, mais l'obscurité dissimulait encore les intrus. Le groupe attendit qu'il s'éloigne, puis se remit en mouvement plus lentement encore.

Destos et Morrÿs surveillèrent si attentivement le sol qu'ils oublièrent les branches à hauteur de visage. L'un d'eux en heurta une de plein fouet et ne parvint pas à retenir un cri. Le ptéranodon en patrouille replia aussitôt ses ailes et fondit vers eux. L'effet de surprise était perdu, mais Azaka leur souffla qu'ils pouvaient encore sauver leur approche s'ils l'abattaient sans alerter les autres.

Morrÿs lança un rayon glacé qui frappa la créature en plein vol et ralentit son approche. Azaka dégaina son arc et planta une flèche dans sa chair avant que son second tir ne se perde dans la nuit. Trop éloigné pour employer ses autres prières, Destos appela sa bénédiction sur Azaka, Kael et lui-même. Obélik attendit, une hachette en main, que leur adversaire entre enfin à portée.

Kael désigna le ptéranodon comme sa proie et banda son arc. Sa flèche toucha grâce à la bénédiction de Destos et arracha la créature à sa trajectoire. Elle perdit brusquement de l'altitude, se rétablit de justesse, puis reprit son plongeon en direction du rôdeur. Morrÿs l'atteignit d'un second rayon glacé et, lorsqu'elle passa suffisamment près, Obélik lança sa hachette et la blessa à son tour.

Le ptéranodon percuta presque Kael en abattant ses griffes sur lui. La violence du choc déchira ses protections et une terreur brutale saisit l'orc. Déjà criblée de projectiles et déséquilibrée, la créature tenta ensuite de le mordre. Dans la confusion, ses mâchoires se refermèrent sur sa propre aile.

Le ptéranodon s'acheva tout seul.

\illustration[0.95\textwidth]{2026-11-10/combat-pterafolk.png}
{\textit{Sous la lumière de Firefinger, les aventuriers tentèrent d'abattre le ptéranodon avant qu'il n'alerte toute la colonie.}}
{combat_pteranodon_firefinger}

Personne ne célébra la victoire. Au-dessus d'eux, le reste de la colonie dormait encore et le moindre cri pouvait transformer cette première rencontre en désastre. Ils reprirent donc leur avancée en contournant le corps. Destos réussit malgré tout à trébucher dessus et s'étala dans un bruit étouffé, sans réveiller les créatures perchées plus haut.

Quelques instants plus tard, les quatre aventuriers et leur guide atteignirent enfin la base de Firefinger. Devant eux, la première échelle se dressait dans l'obscurité jusqu'à une ouverture creusée dans la roche. Le sommet enflammé paraissait encore terriblement loin.
